\documentclass[11pt,a4paper]{article}
\usepackage[margin=2.2cm]{geometry}
\usepackage[T1]{fontenc}
\usepackage[utf8]{inputenc}
\usepackage{enumitem}
\usepackage{booktabs}
\usepackage{tabularx}
\usepackage{longtable}
\usepackage{xcolor}
\usepackage[most]{tcolorbox}
\usepackage{hyperref}

\definecolor{qblue}{RGB}{0,120,180}
\newtcolorbox{question}[1]{colback=qblue!6,colframe=qblue,fonttitle=\bfseries,
  title={Exercise #1},breakable}
\newcommand{\answer}{\par\noindent\textbf{Answer.}\ }
\newcommand{\refsec}[1]{\hfill\textit{\small[Ref.\ Sec.~#1]}}

\newcommand{\spec}[2]{\noindent\begin{tabularx}{\linewidth}{@{}>{\bfseries}p{3.3cm}X@{}}\toprule
\multicolumn{2}{@{}l}{\textbf{#1}}\\\midrule #2 \bottomrule\end{tabularx}\par\medskip}

\title{Software Engineering (10th ed.) --- Chapter 4: Requirements Engineering\\
\large Answers to Exercises 4.1--4.10}
\author{Nguyễn Đình Thiên Lộc}
\date{\today}

\begin{document}
\maketitle
\noindent\textit{Section references (e.g.\ Sec.~4.1.2) point to the sections of Chapter~4 of
I.~Sommerville, \emph{Software Engineering}, 10th ed., where the supporting material is found.}

\tableofcontents
\newpage

%=====================================================================
\section*{Exercise 4.1}\addcontentsline{toc}{section}{Exercise 4.1 -- Types of requirements}
\begin{question}{4.1}
Identify and briefly describe four types of requirements that may be defined for a computer-based system.
\end{question}
\answer
\begin{enumerate}[leftmargin=*]
  \item \textbf{User requirements} -- statements in natural language plus diagrams of the services the
  system provides and the constraints under which it operates, written for customers and end-users.
  \refsec{4 (intro)}
  \item \textbf{System requirements} -- more detailed descriptions of the system's functions, services and
  operational constraints; they define exactly what is to be implemented and may form part of the contract.
  \refsec{4 (intro)}
  \item \textbf{Functional requirements} -- statements of the services the system should provide, how it
  should react to particular inputs and behave in particular situations (and sometimes what it should not do).
  \refsec{4.1.1}
  \item \textbf{Non-functional requirements} -- constraints on the services or the system as a whole, such as
  timing, reliability, security, usability, and constraints imposed by standards. They are classified as
  \emph{product}, \emph{organizational} and \emph{external} requirements. \refsec{4.1.2}
\end{enumerate}

%=====================================================================
\section*{Exercise 4.2}\addcontentsline{toc}{section}{Exercise 4.2 -- Ambiguities in the drone description}
\begin{question}{4.2}
Discover ambiguities or omissions in the following statement of the requirements for part of a drone
system intended for search and recovery:\\[2pt]
\textit{The drone, a quad chopper, will be very useful in search and recovery operations, especially in
remote areas or in extreme weather conditions. It will click high-resolution images. It will fly according
to a path preset by a ground operator, but will be able to avoid obstacles on its own, returning to its
original path whenever possible. The drone will also be able to identify various objects and match them
to the target it is looking for.}
\end{question}
\answer Natural-language requirements suffer from lack of clarity, confusion between functional and
non-functional requirements, and amalgamation of several requirements in one sentence
\refsec{4.4.1}. Problems found:
\begin{enumerate}[leftmargin=*]
  \item \emph{``very useful''} -- a goal, not a verifiable requirement \refsec{4.1.2}.
  \item \emph{``remote areas''} -- maximum operating range from the ground station is not stated; nor is
  the communication link used when out of radio range.
  \item \emph{``extreme weather conditions''} -- undefined: which wind speeds, temperatures, rain/snow levels?
  \item \emph{``high-resolution images''} -- no resolution, frame rate, still vs.\ video, day/night (thermal)
  capability; how often are images taken and where are they stored or transmitted?
  \item \emph{``a path preset by a ground operator''} -- how is it specified (waypoints, altitude)? Can it
  be changed during flight? Can the operator take manual control?
  \item \emph{``avoid obstacles''} -- which obstacles (static, moving), at what minimum distance?
  \emph{``whenever possible''} -- what happens when it is not possible (hover, return home, alert operator)?
  \item \emph{``identify various objects''} -- which object classes? How is the target described to the
  drone? Required accuracy / false-positive rate? What does the drone do after a match (notify, hover,
  mark GPS position)?
  \item \textbf{Omissions}: flight endurance and battery behaviour (low-battery return), loss of
  communication behaviour, maximum altitude/speed, payload (e.g.\ dropping supplies for ``recovery''),
  safety and legal constraints (aviation regulations) \refsec{4.1.2}, and response times.
\end{enumerate}

%=====================================================================
\section*{Exercise 4.3}\addcontentsline{toc}{section}{Exercise 4.3 -- Structured specification}
\begin{question}{4.3}
Rewrite the above description using the structured approach described in this chapter. Resolve the
identified ambiguities in a sensible way.
\end{question}
\answer Using the standard form of a structured specification \refsec{4.4.2}
(function, description, inputs, source, outputs, destination, action, requirements, pre/post-conditions,
side effects), with ambiguities resolved by reasonable assumptions:

\spec{Drone/SR/01 -- Follow preset flight path}{
Function & Follow flight path \\
Description & Fly the quad-chopper along a route defined by the ground operator. \\
Inputs & Ordered list of waypoints (latitude, longitude, altitude 30--120\,m), cruise speed ($\le$ 15\,m/s). \\
Source & Ground control station (GCS) operator. \\
Outputs & Current position, heading, battery level (telemetry at 1\,Hz). \\
Destination & GCS. \\
Action & The drone flies from waypoint to waypoint. The operator may upload a new path or take manual
control at any time. If the GCS link is lost for more than 30\,s, or the battery falls below 25\%, the
drone returns to the launch point. \\
Requires & Valid path uploaded; GPS fix; battery $\ge$ 80\% at take-off. \\
Pre-condition & Pre-flight self-test passed. \\
Post-condition & Drone has landed at launch point or at an operator-specified landing point. \\
Side effects & None. \\
}

\spec{Drone/SR/02 -- Avoid obstacles}{
Function & Obstacle avoidance \\
Description & Detect and avoid static or moving obstacles on the flight path. \\
Inputs & Proximity/vision sensor readings. \\
Source & On-board sensors (range $\ge$ 20\,m). \\
Outputs & Avoidance manoeuvre; obstacle event message. \\
Destination & Flight controller; GCS log. \\
Action & When an obstacle is detected within 10\,m on the path, the drone deviates to keep a clearance
of at least 5\,m, then rejoins the original path at the nearest point beyond the obstacle. If no safe
route is found within 30\,s, it hovers and requests operator instructions. \\
Requires & Sensors operational. \\
Pre-condition & Drone is in autonomous flight. \\
Post-condition & Drone is back on the original path, or hovering awaiting the operator. \\
Side effects & Mission time increases. \\
}

\spec{Drone/SR/03 -- Capture images}{
Function & Capture images \\
Description & Record imagery of the search area. \\
Inputs & Camera and thermal sensor data. \\
Source & On-board camera (minimum 12\,MP stills, 4K/30\,fps video) and thermal camera. \\
Outputs & Geo-tagged images, every 2\,s or on operator command. \\
Destination & On-board storage; low-resolution preview streamed to GCS. \\
Action & Images are tagged with GPS position, altitude and time and stored; a preview is transmitted to the GCS. \\
Requires & At least 2\,GB free storage. \\
Pre-condition & Drone airborne. \\
Post-condition & Images stored. \\
Side effects & None. \\
}

\spec{Drone/SR/04 -- Identify target}{
Function & Identify target objects \\
Description & Recognize objects in images and match them to the search target. \\
Inputs & Captured images; target description (class: person, vehicle, boat; colour; reference photo). \\
Source & Drone/SR/03; GCS operator. \\
Outputs & Match alert with confidence score, image and GPS location. \\
Destination & GCS. \\
Action & Each image is classified; any object of the target class with a match confidence $\ge$ 80\% raises
an alert and the drone hovers over the location until the operator confirms or rejects the match. \\
Requires & Target description uploaded before or during the mission. \\
Pre-condition & Drone/SR/03 active. \\
Post-condition & Alert acknowledged by operator. \\
Side effects & Mission path is paused while waiting. \\
}

\noindent\textbf{Environmental constraint (non-functional):} the drone shall operate in winds up to
40\,km/h, in rain up to 10\,mm/h, and at temperatures from $-20^\circ$C to $+45^\circ$C, at up to 5\,km from the GCS.

%=====================================================================
\section*{Exercise 4.4}\addcontentsline{toc}{section}{Exercise 4.4 -- Non-functional requirements}
\begin{question}{4.4}
Write a set of non-functional requirements for the drone system, setting out its expected safety and
response time.
\end{question}
\answer Non-functional requirements should be written quantitatively so they can be tested, using
metrics such as speed, reliability or robustness \refsec{4.1.2}.

\smallskip\noindent\textbf{Safety (product requirements)}
\begin{enumerate}[label=NFR-S\arabic*,leftmargin=*]
  \item If the communication link with the GCS is lost for more than 30\,s, the drone shall automatically return to its launch point.
  \item When battery charge falls below 25\%, the drone shall abort the mission and return; below 10\% it shall perform an immediate controlled landing.
  \item The drone shall maintain a minimum separation of 5\,m from any detected obstacle and 30\,m from people not designated as targets.
  \item The drone shall not exceed an altitude of 120\,m above ground level nor enter geo-fenced no-fly zones (external/legislative requirement).
  \item A manual operator override command shall take precedence over all autonomous behaviour.
  \item The probability of an uncontrolled crash shall be less than 1 per 10{,}000 flight hours.
\end{enumerate}

\noindent\textbf{Response time (performance)}
\begin{enumerate}[label=NFR-R\arabic*,leftmargin=*]
  \item The drone shall start an avoidance manoeuvre within 200\,ms of detecting an obstacle.
  \item Operator commands shall be executed within 500\,ms of being sent from the GCS.
  \item A target-match alert shall reach the GCS within 2\,s of the image being captured.
  \item The live video preview latency at the GCS shall not exceed 1\,s.
  \item Telemetry (position, battery) shall be updated at the GCS at least once per second.
  \item The drone shall be ready for take-off within 3\,minutes of power-on.
\end{enumerate}

%=====================================================================
\section*{Exercise 4.5}\addcontentsline{toc}{section}{Exercise 4.5 -- Standard-format user requirements}
\begin{question}{4.5}
Using the technique suggested here, where natural language descriptions are presented in a standard
format, write plausible user requirements for the following functions:
\begin{itemize}
  \item An unattended petrol (gas) pump system that includes a credit card reader. The customer swipes the
  card through the reader, then specifies the amount of fuel required. The fuel is delivered and the
  customer's account debited.
  \item The cash-dispensing function in a bank ATM.
  \item In an Internet banking system, a facility that allows customers to transfer funds from one account
  held with the bank to another account with the same bank.
\end{itemize}
\end{question}
\answer Guidelines followed \refsec{4.4.1}: a standard format, ``shall'' for mandatory and ``should''
for desirable requirements, one requirement per statement, plus a rationale.

\subsection*{(a) Unattended petrol pump}
\begin{enumerate}[label=PUMP-\arabic*,leftmargin=*]
  \item The system \textbf{shall} allow the customer to pay by swiping a credit or debit card in the card reader.\\
  \textit{Rationale:} the pump is unattended, so no cashier is available.
  \item The system \textbf{shall} validate the card with the card issuer and reject expired, blocked or unreadable cards with an explanatory message.
  \item After card validation the system \textbf{shall} ask the customer to enter the amount of fuel, either as a volume (litres) or a monetary value, up to a pre-authorized maximum.
  \item The system \textbf{shall} pre-authorize the requested amount on the card account before enabling the pump.\\
  \textit{Rationale:} prevents delivery of fuel that cannot be paid for.
  \item The system \textbf{shall} stop delivery automatically when the requested amount is reached or the nozzle is replaced.
  \item The system \textbf{shall} debit the customer's account with the cost of the fuel actually delivered, not the requested amount.
  \item The system \textbf{should} offer a printed or e-mailed receipt.
\end{enumerate}

\subsection*{(b) ATM cash dispensing}
\begin{enumerate}[label=ATM-\arabic*,leftmargin=*]
  \item The ATM \textbf{shall} dispense cash only after the customer has inserted a valid card and entered the correct PIN.
  \item The ATM \textbf{shall} allow the customer to select a predefined amount or enter an amount that is a multiple of the smallest note available.
  \item The ATM \textbf{shall} refuse a withdrawal that exceeds the account balance/overdraft limit or the daily withdrawal limit, and inform the customer of the reason.
  \item The ATM \textbf{shall} return the card before dispensing the cash.\\
  \textit{Rationale:} reduces the number of cards left in the machine.
  \item The ATM \textbf{shall} debit the account only after the cash has been successfully dispensed; if dispensing fails, no debit shall occur.
  \item The ATM \textbf{shall} retract cash that is not taken within 30\,s and record the event.
  \item The ATM \textbf{should} offer a receipt showing the amount withdrawn and remaining balance.
\end{enumerate}

\subsection*{(c) Internet banking -- internal funds transfer}
\begin{enumerate}[label=TRF-\arabic*,leftmargin=*]
  \item The system \textbf{shall} allow an authenticated customer to transfer funds from one of their accounts to any account held at the same bank.
  \item The customer \textbf{shall} specify the source account, destination account number, amount and an optional reference.
  \item The system \textbf{shall} display the name of the destination account holder and ask the customer to confirm before executing the transfer.\\
  \textit{Rationale:} reduces transfers to wrong accounts.
  \item The system \textbf{shall} reject transfers that exceed the available balance or the customer's daily transfer limit.
  \item Transfers above a bank-defined threshold \textbf{shall} require a one-time password as second-factor authentication.
  \item The transfer \textbf{shall} be applied to both accounts atomically and immediately, and a confirmation with a transaction ID shown.
  \item The system \textbf{should} allow the customer to save frequent destination accounts as payees.
\end{enumerate}

%=====================================================================
\section*{Exercise 4.6}\addcontentsline{toc}{section}{Exercise 4.6 -- Tracking FR/NFR relationships}
\begin{question}{4.6}
Suggest how an engineer responsible for drawing up a system requirements specification might keep track
of the relationships between functional and non-functional requirements.
\end{question}
\answer Non-functional requirements often affect many functional requirements, and one NFR can generate
several functional requirements \refsec{4.1.2}. To keep track:
\begin{enumerate}[leftmargin=*]
  \item \textbf{Unique identification} -- give every requirement a unique ID (e.g.\ FR-12, NFR-PERF-3) so it can be cross-referenced \refsec{4.6.1}.
  \item \textbf{Traceability matrix} -- build a table with functional requirements as rows and non-functional requirements as columns; mark each cell where the NFR constrains the FR (``D'' for depends, ``R'' for weaker relation). This makes impact analysis easy when either changes \refsec{4.6.1, 4.6.2}.
  \item \textbf{Explicit links in the specification} -- in each structured specification form \refsec{4.4.2}, include a field ``Related NFRs'', and for each NFR list the FRs it constrains or was derived from.
  \item \textbf{Record rationale and source} -- note why each NFR exists and which stakeholder proposed it, so conflicts can be traced to their origin.
  \item \textbf{Tool support} -- use a requirements management tool (e.g.\ IBM DOORS) or at least a shared spreadsheet/database that stores requirements and links and can report affected requirements automatically \refsec{4.6.1}.
\end{enumerate}

%=====================================================================
\section*{Exercise 4.7}\addcontentsline{toc}{section}{Exercise 4.7 -- ATM use cases}
\begin{question}{4.7}
Using your knowledge of how an ATM is used, develop a set of use cases that could serve as a basis for
understanding the requirements for an ATM system.
\end{question}
\answer Use cases identify the actors in an interaction and name the type of interaction; each is then
described in more detail \refsec{4.4.3}. Scenarios can be written for each \refsec{4.3.2}.

\noindent\textbf{Actors:} Bank customer, Bank system (account database), ATM operator/maintenance staff, Card issuer network.

\begin{longtable}{@{}p{3.2cm}p{2.8cm}p{9.2cm}@{}}
\toprule
\textbf{Use case} & \textbf{Actors} & \textbf{Description} \\ \midrule\endhead
Validate user & Customer, Bank & Customer inserts card and enters PIN; system checks card validity and PIN; card retained after 3 failed attempts. Included in all customer use cases. \\
Withdraw cash & Customer, Bank & Customer selects amount; bank checks balance and limits; card returned; cash dispensed; account debited; optional receipt. \\
Check balance & Customer, Bank & Displays or prints current and available balance. \\
Deposit cash/cheque & Customer, Bank & Customer inserts notes or cheque; ATM counts/scans; amount confirmed; account credited (cheques pending clearance). \\
Transfer funds & Customer, Bank & Transfer between the customer's linked accounts. \\
Change PIN & Customer, Bank & Customer enters old PIN and new PIN twice; bank updates PIN. \\
Print mini-statement & Customer, Bank & Prints the last 10 transactions. \\
Use foreign card & Customer, Card network & Card of another bank is validated through the inter-bank network; fee shown before confirmation. \\
Replenish ATM & Operator & Operator authenticates, reloads cash cassettes and paper, records the amounts loaded. \\
Run maintenance & Operator & Diagnostics, clear jams, retrieve retained cards and deposits. \\
Report status & ATM, Bank & ATM periodically reports cash level and faults to the bank. \\
\bottomrule
\end{longtable}

\noindent\textbf{Example detailed use case -- Withdraw cash}
\begin{description}[leftmargin=2.6cm,style=nextline]
  \item[Actors] Customer, Bank system.
  \item[Pre-condition] Customer validated (\emph{Validate user}).
  \item[Normal flow] 1. Customer chooses ``Withdraw''. 2. Selects account and amount. 3. ATM requests authorization from bank. 4. Bank approves. 5. Card returned. 6. Cash dispensed. 7. Account debited; receipt printed if requested.
  \item[Alternatives] Insufficient funds or limit exceeded $\rightarrow$ message, return to menu. ATM lacks notes $\rightarrow$ suggest a different amount. Cash not taken in 30\,s $\rightarrow$ retracted, transaction reversed.
  \item[Post-condition] Account balance reduced by the dispensed amount; transaction logged.
\end{description}

%=====================================================================
\section*{Exercise 4.8}\addcontentsline{toc}{section}{Exercise 4.8 -- Two scribes in a review}
\begin{question}{4.8}
To minimize mistakes during a requirements review, an organization decides to allocate two scribes to
document the review session. Explain how this can be done.
\end{question}
\answer In a requirements review a team of customer and developer staff reads the document and checks
for errors, anomalies and inconsistencies, which must be recorded before being negotiated \refsec{4.5}.
Two scribes can be used as follows:
\begin{enumerate}[leftmargin=*]
  \item \textbf{Divide responsibilities.} Scribe A records the \emph{problems} raised (requirement ID,
  description, type -- ambiguity, inconsistency, omission, unrealistic, unverifiable -- and who raised it);
  Scribe B records the \emph{discussion and decisions} (agreed resolution, actions, owner, deadline).
  Alternatively, they may split by document section.
  \item \textbf{Use a common template} -- both use the same requirement IDs and problem categories
  (validity, consistency, completeness, realism, verifiability \refsec{4.5}) so records can be merged.
  \item \textbf{Independent overlapping notes} -- for important items both scribes record independently;
  differences highlight possible misunderstandings.
  \item \textbf{Reconcile after the session} -- the scribes compare and merge their notes into one review
  report soon after the meeting, resolving discrepancies (with the moderator if needed).
  \item \textbf{Confirm with participants} -- the consolidated report is circulated to reviewers, customer
  and developers for confirmation; it then feeds into the change management process \refsec{4.6.2}.
\end{enumerate}

%=====================================================================
\section*{Exercise 4.9}\addcontentsline{toc}{section}{Exercise 4.9 -- Emergency change process}
\begin{question}{4.9}
When emergency changes have to be made to systems, the system software may have to be modified before
changes to the requirements have been approved. Suggest a model of a process for making these
modifications that will ensure that the requirements document and the system implementation do not
become inconsistent.
\end{question}
\answer The normal change process is \emph{problem analysis and change specification $\rightarrow$
change analysis and costing $\rightarrow$ change implementation} \refsec{4.6.2}. If urgent changes are
made to the system first, it is easy to forget to update the requirements document, so it must be
updated as soon as possible \refsec{4.6.2}. A suggested emergency process:
\begin{enumerate}[leftmargin=*]
  \item \textbf{Raise an emergency change request (ECR)} -- record the problem, the affected
  requirement IDs and the reason the change is urgent. The ECR is logged in the change management tool
  with status \emph{``Emergency -- requirements pending''}.
  \item \textbf{Fast authorization} -- a small authority (e.g.\ project manager plus customer
  representative) approves the emergency fix, instead of the full change board.
  \item \textbf{Implement the fix} -- code is modified; every commit references the ECR ID, and a brief
  draft of the changed requirement is attached to the ECR.
  \item \textbf{Mandatory follow-up} -- the ECR cannot be closed until the requirements document is
  updated. The system is labelled as a ``temporary'' release.
  \item \textbf{Retrospective change analysis} -- within a fixed deadline (e.g.\ 5 working days) the
  change goes through normal analysis and costing using traceability links \refsec{4.6.1} to find other
  affected requirements and design.
  \item \textbf{Approve and update} -- the change board formally approves (or requests rework of) the
  change; the requirements document, design and tests are updated, and the ECR is closed.
  \item \textbf{Audit} -- periodic reports list all open ECRs, so no emergency change remains undocumented.
\end{enumerate}

%=====================================================================
\section*{Exercise 4.10}\addcontentsline{toc}{section}{Exercise 4.10 -- Ethical dilemma}
\begin{question}{4.10}
You have taken a job with a software user who has contracted your previous employer to develop a system
for them. You discover that your company's interpretation of the requirements is different from the
interpretation taken by your previous employer. Discuss what you should do in such a situation. You know
that the costs to your current employer will increase if the ambiguities are not resolved. However, you
also have a responsibility of confidentiality to your previous employer.
\end{question}
\answer This is a professional-ethics problem. Requirements are often ambiguous and different
stakeholders interpret them differently \refsec{4.4.1}; resolving that is the purpose of requirements
validation and negotiation \refsec{4.3, 4.5}. A reasonable course of action:
\begin{enumerate}[leftmargin=*]
  \item \textbf{Do not disclose confidential information} of the previous employer (internal
  designs, documents, costings, strategies). Confidentiality obligations continue after leaving.
  \item \textbf{Raise the issue using only public/shared information.} The requirements document is
  shared between the customer (your current employer) and the contractor, so you can legitimately point
  out that particular requirements are \emph{ambiguous} and could be interpreted in more than one way,
  without revealing how the previous employer interprets them.
  \item \textbf{Recommend a formal requirements review} \refsec{4.5} with the contractor, in which both
  parties state their interpretation of the ambiguous requirements; the differences will then appear
  naturally and can be negotiated \refsec{4.3.1}.
  \item \textbf{Declare the conflict of interest} to your current manager and, if necessary, ask not to be
  involved in negotiations where your previous knowledge would be relevant.
  \item \textbf{If unsure, seek advice} (e.g.\ professional code of ethics, legal department) before
  acting.
\end{enumerate}
In this way you protect your current employer's interests (the ambiguity gets resolved early, when it is
cheapest) while honouring your duty of confidentiality to the previous employer.

\section*{Reference}
\addcontentsline{toc}{section}{Reference}
I.~Sommerville, \emph{Software Engineering}, 10th ed., Pearson, 2016, Chapter~4 ``Requirements engineering''
(Sec.~4.1 Functional and non-functional requirements; 4.2 RE processes; 4.3 Requirements elicitation;
4.4 Requirements specification; 4.5 Requirements validation; 4.6 Requirements change).

\end{document}
